% Consolidated 10-thermo-first-law from 7c59eacf2bbc4f2718617404cd46e24a47721d7d; original TeX fragments preserved.

% === BEGIN 10-thermo-first-law-026-035.tex blob 4ca2c62a4e5397a1dae95c40c70c1d7f27190de5 ===
% Irodov, Part II "Thermodynamics and Molecular Physics",
% section 2.2 "The First Law of Thermodynamics. Heat Capacity";
% original numbered problems 26--35 (printed pp.65--66 / PDF pp.74--75).
% Gas molar masses M and heat capacities are ideal-gas approximations.

\thermoenergyproblem{26}{65}
证明：若室内空气与外界相通、外界压强
$p$ 保持恒定，则\emph{室内空气的总内能}
与温度无关。
若室内容积 $V=40\,\mathrm{m^3}$，
空气压强为标准大气压，
估计室内空气的内能 $U$。
\exsolution
将室内空气视为双原子理想气体，
暂不激发振动自由度，
其绝热指数 $\gamma=C_p/C_V\simeq7/5$。
室内气体总内能为
\[
U=\nu C_VT,
\]
其中 $\nu$ 是室内气体的摩尔数。
因为理想气体状态方程
$pV=\nu RT$，且
$C_V=R/(\gamma-1)$，
故
\[
\boxed{U=\frac{pV}{\gamma-1}}.
\]
当房间容积 $V$ 和外界压强 $p$ 固定，
$U$ 不随温度变化；升温时部分空气流出，
使摩尔数 $\nu\propto1/T$。
代入 $\gamma-1=0.4$、
$p=1.01325\times10^5\,\mathrm{Pa}$、
$V=40\,\mathrm{m^3}$：
\[
\boxed{U\simeq\frac{1.01325\times10^5\cdot40}{0.4}
\simeq1.0\times10^7\,\mathrm J=10\,\mathrm{MJ}}.
\]
\excheck 这一结论要求\emph{房间与外界通气}；
密闭、质量不变的房间内理想气体内能仍随温度变化。

\thermoenergyproblem{27}{66}
一个绝热封闭容器中装有理想气体，
其摩尔质量为 $M$、
热容比为 $\gamma=C_p/C_V$。
容器及气体原先整体以速率 $v$
平动，突然停止。
若停止过程的气体宏观平动动能
最终完全转化为其内能，求气体升温 $\Delta T$。
\exsolution
设气体总质量为 $m$，摩尔数
$\nu=m/M$。
停止前相对实验室的气体宏观平动动能为
$K=\frac12mv^2$。
绝热且假定没有把能量带到容器外，
这部分宏观动能转化为热内能：
\[
\nu C_V\Delta T=\frac12mv^2.
\]
由
$C_V=R/(\gamma-1)$、
$\nu=m/M$，消去质量：
\[
\frac mM\frac{R}{\gamma-1}\Delta T
=\frac12mv^2,
\]
得到
\[
\boxed{\Delta T=\frac{(\gamma-1)Mv^2}{2R}}.
\]
\excheck 这不是任意急停过程的无条件能量守恒结论；
假定气体原有宏观平动能全部内化，
且容器壁、外部制动机构不吸收相应份额。

\thermoenergyproblem{28}{66}
两只绝热刚性容器由带阀门的短管相连，
内部装的是同一种理想气体（空气）。
两容器容积、初始压强和温度分别为
$(V_1,p_1,T_1)$ 与 $(V_2,p_2,T_2)$。
打开阀门，最终达到共同平衡温度和压强。
求末温 $T$ 与末压 $p$。
\exsolution
设两容器最初摩尔数为
\[
\nu_i=\frac{p_iV_i}{RT_i}\quad(i=1,2).
\]
整个系统刚性、绝热，外界没有机械功
与热量交换，因此总内能守恒：
\[
(\nu_1+\nu_2)C_VT
=\nu_1C_VT_1+\nu_2C_VT_2.
\]
消去同一种气体的 $C_V$，
\[
T=\frac{\nu_1T_1+\nu_2T_2}{\nu_1+\nu_2}
=\boxed{\frac{p_1V_1+p_2V_2}
{p_1V_1/T_1+p_2V_2/T_2}}
=\boxed{\frac{T_1T_2(p_1V_1+p_2V_2)}
{p_1V_1T_2+p_2V_2T_1}}.
\]
总摩尔数守恒且总体积为 $V_1+V_2$，
于是
\[
p(V_1+V_2)=(\nu_1+\nu_2)RT
=p_1V_1+p_2V_2,
\]
即
\[
\boxed{p=\frac{p_1V_1+p_2V_2}{V_1+V_2}}.
\]
\excheck 平衡温度是初始\emph{摩尔数加权平均}，
不简单等于 $(T_1+T_2)/2$。

\thermoenergyproblem{29}{66}
密闭容器装有初始标准状况的氢气，
容积为 $V=5.0\,\mathrm L$，
将其冷却
$\Delta T=55\,\mathrm K$。
求气体内能变化以及气体放出的热量。
\exsolution
标准状况取
$T_0\simeq273.15\,\mathrm K$、
$p_0\simeq1.013\times10^5\,\mathrm{Pa}$。
理想氢气在此温度视为双原子气体，
$C_V=5R/2=R/(\gamma-1)$，
其中 $\gamma\simeq7/5$。
密闭刚性容器体积不变，
因此气体没有对外体积功。
初始摩尔数
$\nu=p_0V/(RT_0)$，
降温的内能变化
\[
\Delta U=\nu C_V(-\Delta T)
=-\frac{p_0V\Delta T}{T_0(\gamma-1)}.
\]
代入 $V=5.0\times10^{-3}\,\mathrm{m^3}$：
\[
\boxed{\Delta U\simeq-2.5\times10^2\,\mathrm J
=-0.25\,\mathrm{kJ}}.
\]
热力学第一定律
$Q=\Delta U+A$，等容过程 $A=0$，
故气体\emph{吸收}的热为 $Q=\Delta U<0$，
向环境放出的热量大小为
\[
\boxed{Q_{\rm out}=-Q=-\Delta U\simeq0.25\,\mathrm{kJ}}.
\]
\excheck 区分带符号的气体吸热 $Q$ 与
放热量的正值 $Q_{\rm out}$。

\thermoenergyproblem{30}{66}
对氮气进行等压加热，
使气体在过程中对外做功
$A=2.0\,\mathrm J$。
假设氮气为常温双原子理想气体，
求需提供的热量 $Q$。
\exsolution
等压理想气体体积功
$A=\nu R\Delta T$。
吸热为
$Q=\nu C_p\Delta T$。
相除：
\[
\frac QA=\frac{C_p}{R}.
\]
由 $C_p-C_V=R$ 和
$\gamma=C_p/C_V$，
得 $C_p=\gamma R/(\gamma-1)$。
因此
\[
\boxed{Q=\frac{\gamma}{\gamma-1}A}.
\]
氮气 $\gamma\simeq7/5$，
故
\[
\boxed{Q=\frac{1.4}{0.4}(2.0\,\mathrm J)
\simeq7.0\,\mathrm J}.
\]
\excheck 供热不仅成为外功，还使气体内能增加；
所以 $Q>A$。

\thermoenergyproblem{31}{66}
一摩尔某理想气体经等压加热，
温度上升 $\Delta T=72\,\mathrm K$，
吸热 $Q=1.60\,\mathrm{kJ}$。
求气体做功 $A$、内能增量 $\Delta U$
和热容比 $\gamma=C_p/C_V$。
\exsolution
等压理想气体做功
\[
\boxed{A=R\Delta T
=(8.314)(72)\,\mathrm J\simeq0.60\,\mathrm{kJ}}.
\]
由第一定律
$Q=\Delta U+A$，
\[
\boxed{\Delta U=Q-A\simeq1.00\,\mathrm{kJ}}.
\]
定压摩尔热容
$C_p=Q/\Delta T$，
定容摩尔热容
$C_V=(Q-R\Delta T)/\Delta T$，
故
\[
\boxed{\gamma=\frac{C_p}{C_V}
=\frac{Q}{Q-R\Delta T}
\simeq1.60}.
\]
\excheck 用 $\gamma\simeq1.60$ 反算
$C_V=R/(\gamma-1)\simeq13.9\,\mathrm{J/(mol\,K)}$，
与 $\Delta U/\Delta T$ 一致。

\thermoenergyproblem{32}{66}
某理想气体共 $\nu=2.0$ mol，
初温 $T_0=300\,\mathrm K$。
先等容冷却，使压强降低
$n=2.0$ 倍；
再等压膨胀，使温度恢复到初始值 $T_0$。
求整个过程从外界\emph{净吸收}的热量。
\exsolution
等容第一段由 $p\propto T$，
末温为 $T_1=T_0/n$。
第二段等压加热使温度
从 $T_1$ 回到 $T_0$。
两段吸热总和为
\[
Q=\nu C_V(T_1-T_0)
+\nu C_p(T_0-T_1).
\]
由 $C_p-C_V=R$，得
\[
Q=\nu R(T_0-T_1)
=\nu RT_0\left(1-\frac1n\right).
\]
因此
\[
\boxed{Q=(2.0)(8.314)(300)(1-1/2)
\simeq2.5\,\mathrm{kJ}}.
\]
\excheck 第一段放热、第二段吸热，
\emph{净吸热}为正，且不依赖气体热容比。

\thermoenergyproblem{33}{66}
某理想混合气由
$\nu_1=2.0$ mol 氧气、
$\nu_2=3.0$ mol 二氧化碳组成。
求混合气的热容比 $\gamma$，
以两气体各自热容比
$\gamma_1,\gamma_2$ 表示并给出估算值。
\exsolution
混合气的总定容热容为各组分热容之和：
\[
C_{V,\Sigma}=
\nu_1\frac R{\gamma_1-1}
+\nu_2\frac R{\gamma_2-1}.
\]
总定压热容同样相加：
\[
C_{p,\Sigma}=
\nu_1\frac{\gamma_1 R}{\gamma_1-1}
+\nu_2\frac{\gamma_2 R}{\gamma_2-1}.
\]
所以混合气的热容比为
\[
\boxed{\gamma=
\frac{\nu_1\gamma_1(\gamma_2-1)
+\nu_2\gamma_2(\gamma_1-1)}
{\nu_1(\gamma_2-1)+\nu_2(\gamma_1-1)}}.
\]
如果使用原书温区下的常见近似
$\gamma_1\simeq1.40$（氧气）、
$\gamma_2\simeq1.30$（二氧化碳），
\[
\boxed{\gamma\simeq
\frac{2(1.40)(0.30)+3(1.30)(0.40)}
{2(0.30)+3(0.40)}
=\frac{2.40}{1.80}\simeq1.33}.
\]
\excheck CO$_2$ 的有效热容比依赖温度和
是否激发振动自由度；题目未给具体温度，
数值结果需说明所采用的组分热容假设。

\thermoenergyproblem{34}{66}
由质量为 $7.0\,\mathrm g$ 的氮气和
$20\,\mathrm g$ 的氩气组成理想混合气。
在常温、热容近似不随温度变化时，
求混合气的质量比热容
$c_V,c_p$。
\exsolution
氮气为双原子气体，摩尔热容约
$C_{V,{\rm N_2}}=5R/2$、
$C_{p,{\rm N_2}}=7R/2$；
氩气为单原子气体，分别为
$3R/2$ 与 $5R/2$。
两种气体的摩尔数：
\[
\nu_{\rm N_2}=\frac{7.0}{28}=0.25\,\mathrm{mol},
\qquad
\nu_{\rm Ar}=\frac{20}{40}=0.50\,\mathrm{mol}.
\]
混合气总质量 $m_{\Sigma}=27.0\,\mathrm g$。
以每克质量计的定容比热为
\[
c_V=
\frac{\nu_{\rm N_2}(5R/2)+\nu_{\rm Ar}(3R/2)}
{m_{\Sigma}}
=\frac{(0.625+0.75)R}{27.0\,\mathrm g},
\]
故
\[
\boxed{c_V\simeq0.42\,\mathrm{J/(g\,K)}}.
\]
同样
\[
c_p=
\frac{\nu_{\rm N_2}(7R/2)+\nu_{\rm Ar}(5R/2)}
{m_{\Sigma}}
=\frac{(0.875+1.25)R}{27.0\,\mathrm g},
\]
所以
\[
\boxed{c_p\simeq0.655\,\mathrm{J/(g\,K)}}.
\]
\excheck 组分按\emph{摩尔数乘摩尔热容}求总热容，
最后除以总质量；不能直接对质量比热做无权平均。

\thermoenergyproblem{35}{66}
一摩尔理想气体在竖直圆筒中，
被无质量、无摩擦活塞封住，
活塞上方连通大气。
初温为 $T$，
现缓慢向上提起活塞，
使气体\emph{等温}膨胀到原体积的
$n$ 倍。
求外界提起活塞需要做的功 $A_{\rm ext}$。
\exsolution
气体初始处于平衡，
因此初始气压
$p_{\rm gas}(V_0)=p_0$。
一摩尔理想气体
$p_{\rm gas}(V)=RT/V$，
且初态 $p_0V_0=RT$。
缓慢提拉活塞的外力大小为
$F_{\rm ext}=[p_0-p_{\rm gas}(V)]S$，
其中 $S$ 为活塞横截面积，
允许外力由向上逐渐变大。
若活塞上升 $h$，体积变化
$\mathrm dV=S\,\mathrm dh$，
因此外力所做功
\[
A_{\rm ext}
=\int_{V_0}^{nV_0}
\left(p_0-\frac{RT}{V}\right)\mathrm dV
=p_0V_0(n-1)-RT\ln n.
\]
由 $p_0V_0=RT$，
\[
\boxed{A_{\rm ext}=RT(n-1-\ln n)}.
\]
\excheck 对 $n\ge1$，
不等式 $n-1\ge\ln n$ 保证外力功非负；
它不同于气体自身的等温膨胀功 $RT\ln n$。
% === END 10-thermo-first-law-026-035.tex ===

% === BEGIN 10-thermo-first-law-036-045.tex blob 2dd37e7687985b55da702755f80bbccb2f093c0d ===
% Irodov, Part II Section 2.2, original numbered problems 36--45,
% printed pp.66--67 (PDF pp.75--76). Fig. 13 in original Answers
% gives qualitative P-T and V-T diagrams, here reproduced in TikZ.

\thermoenergyproblem{36}{66}
一根两端封闭的水平圆筒被一个可无摩擦移动的
活塞分成体积均为 $V_0$ 的左右两腔。
两侧各装有压强 $p_0$、温度相同的理想气体，
初始处于平衡。
现缓慢移动活塞，使其中一腔体积变成另一腔的
$\eta$ 倍，整个过程中两腔气体\emph{等温}。
求外界移动活塞需做的功。
\exsolution
整个圆筒体积保持
$V_1+V_2=2V_0$。
题给最终体积比 $V_1/V_2=\eta$，
因此
\[
V_1=\frac{2\eta}{1+\eta}V_0,\qquad
V_2=\frac2{1+\eta}V_0.
\]
初始时两腔各自的理想气体摩尔数满足
$\nu RT=p_0V_0$。
两腔等温膨胀功之和为
\[
W_{\rm gas}=p_0V_0\ln\frac{V_1}{V_0}
+p_0V_0\ln\frac{V_2}{V_0}
=p_0V_0\ln\frac{4\eta}{(1+\eta)^2}.
\]
因 $(1+\eta)^2\ge4\eta$，
$W_{\rm gas}\le0$；
外力的准静态做功为其相反数：
\[
\boxed{A_{\rm ext}=p_0V_0
\ln\frac{(1+\eta)^2}{4\eta}}.
\]
\excheck 对 $\eta=1$，活塞不发生净位移且
$A_{\rm ext}=0$；对偏离1的比例，外力功非负。

\thermoenergyproblem{37}{67}
三摩尔理想气体初温
$T_0=273\,\mathrm K$。
先等温膨胀至初始体积的
$n=5.0$ 倍，再等容加热，
使最终压强回到最初数值。
两过程总吸热量为
$Q=80\,\mathrm{kJ}$。
求气体的热容比
$\gamma=C_p/C_V$。
\exsolution
第一段等温，
$V_1=nV_0$、$T_1=T_0$，
吸热等于气体做功：
\[
Q_1=\nu RT_0\ln n,\qquad \nu=3.
\]
第一段结束时压强为 $p_0/n$。
第二段等容使压强回到 $p_0$，
故末温为 $T_2=nT_0$。
其吸热
\[
Q_2=\nu C_V(T_2-T_1)
=\nu\frac R{\gamma-1}(n-1)T_0.
\]
总吸热为
$Q=\nu RT_0[\ln n+(n-1)/(\gamma-1)]$，
故
\[
\boxed{\gamma
=1+\frac{n-1}{Q/(\nu RT_0)-\ln n}}.
\]
代入
$\nu=3,T_0=273\,\mathrm K,n=5,Q=80\,\mathrm{kJ}$：
\[
\boxed{\gamma\simeq1.4}.
\]
\excheck 初段为等温，\emph{不}改变气体内能；
最终压力回到初压的第二段是等容加热，温度必须升至
$nT_0$。

\thermoenergyproblem{38}{67}
定性画出理想气体的
等容、等压、等温、绝热四种准静态过程曲线，
分别采用坐标
(a) $(T,p)$；
(b) $(T,V)$，其中 $T$ 为横轴。
\exsolution
从理想气体状态方程
$pV=\nu RT$ 和可逆绝热方程
$pV^\gamma=\text{常数}$ 出发：
\[
\begin{array}{c|c|c}
\text{过程}&p\text{ 对 }T&V\text{ 对 }T\\ \hline
\text{等容}&p\propto T&V=\text{常数}\\
\text{等压}&p=\text{常数}&V\propto T\\
\text{等温}&T=\text{常数}&T=\text{常数}\\
\text{绝热}&p\propto T^{\gamma/(\gamma-1)}
&V\propto T^{-1/(\gamma-1)}
\end{array}
\]
因此在 $(T,p)$ 坐标中，
等压是水平线、等容是过原点的直线、
等温是竖线、绝热为向上弯曲曲线；
在 $(T,V)$ 坐标中，
等容是水平线、等压是过原点的直线、
等温是竖线、绝热为单调下降的曲线。
以下给出两组示意图，均不表示精确比例：
\begin{center}
\begin{tikzpicture}[scale=0.78,>=Stealth]
\draw[->] (0,0)--(4.3,0) node[right] {$T$};
\draw[->] (0,0)--(0,3.2) node[above] {$p$};
\draw[thick] (0,1.2)--(4.1,1.2) node[right] {\scriptsize 等压};
\draw[thick] (0,0)--(2.8,2.8) node[right] {\scriptsize 等容};
\draw[thick] (2.1,0)--(2.1,2.95) node[above] {\scriptsize 等温};
\draw[thick,domain=0.6:2.8,samples=45]
plot (\x,{0.16*(\x)^2.6});
\node at (3.55,2.55) {\scriptsize 绝热};
\begin{scope}[xshift=7cm]
\draw[->] (0,0)--(4.3,0) node[right] {$T$};
\draw[->] (0,0)--(0,3.2) node[above] {$V$};
\draw[thick] (0,1.4)--(4.0,1.4) node[right] {\scriptsize 等容};
\draw[thick] (0,0)--(3.2,2.65) node[right] {\scriptsize 等压};
\draw[thick] (2.2,0)--(2.2,2.95) node[above] {\scriptsize 等温};
\draw[thick,domain=0.85:3.6,samples=45]
plot (\x,{2.6/(\x)^1.35});
\node at (0.9,2.9) {\scriptsize 绝热};
\end{scope}
\end{tikzpicture}
\end{center}
\excheck 绝热曲线的指数由 $\gamma>1$ 决定；
图示的坐标轴不采用一般教材常见的 $(V,p)$，
而是原题明确要求的 $(T,p)$ 和 $(T,V)$。

\thermoenergyproblem{39}{67}
一摩尔氧气初始温度
$T_0=290\,\mathrm K$，
进行可逆绝热压缩，
压强增加 $\eta=10.0$ 倍。
求：(a) 压缩后的气体温度；
(b) 外界对气体做的功。
\exsolution
取常温氧气为理想双原子气体，
绝热指数 $\gamma\simeq7/5$，
定容摩尔热容
$C_V=R/(\gamma-1)$。
可逆绝热关系
$TV^{\gamma-1}=\text{常数}$、
$pV^\gamma=\text{常数}$，
消去体积：
\[
\frac T{T_0}=
\left(\frac p{p_0}\right)^{(\gamma-1)/\gamma}.
\]
因此
\[
\boxed{T=T_0\eta^{(\gamma-1)/\gamma}
=290\cdot10^{2/7}
\simeq0.56\,\mathrm{kK}}.
\]
绝热条件 $Q=0$，
外界对气体做功
$A_{\rm on}=\Delta U=C_V(T-T_0)$，
于是
\[
\boxed{A_{\rm on}
=\frac{RT_0}{\gamma-1}
\left[\eta^{(\gamma-1)/\gamma}-1\right]
\simeq5.6\,\mathrm{kJ}}.
\]
\excheck 压缩升温且外力做正功；
气体\emph{自身}做功为负 $-A_{\rm on}$，
不能与外力做功混同。

\thermoenergyproblem{40}{67}
某质量的氮气从相同初态开始，
分别以可逆绝热方式和等温方式
把气体体积压缩 $\eta=5.0$ 倍。
求两种压缩中外力所做功的比值
$A_{\rm ad}/A_{\rm iso}$。
\exsolution
设初温 $T_0$、摩尔数 $\nu$、
绝热指数 $\gamma$。
对于绝热压缩，
末温
$T_{\rm ad}=T_0\eta^{\gamma-1}$，
外力做功等于内能增加：
\[
A_{\rm ad}=
\nu C_V(T_{\rm ad}-T_0)
=\frac{\nu RT_0}{\gamma-1}
(\eta^{\gamma-1}-1).
\]
等温压缩外力做功为
\[
A_{\rm iso}=\nu RT_0\ln\eta.
\]
相除：
\[
\boxed{\frac{A_{\rm ad}}{A_{\rm iso}}
=\frac{\eta^{\gamma-1}-1}
{(\gamma-1)\ln\eta}}.
\]
氮气 $\gamma\simeq1.40$，
$\eta=5$，得
\[
\boxed{\frac{A_{\rm ad}}{A_{\rm iso}}
\simeq1.4}.
\]
\excheck 绝热压缩过程中气体温度升高，
因此需要的外力功比等温压缩大。

\thermoenergyproblem{41}{67}
一封闭绝热圆筒内有理想气体，
中间有一个\emph{导热且可无摩擦移动}的活塞。
初始两腔体积相等，
两侧气体温度同为 $T_0$，
随后缓慢推动活塞，
使两腔体积比
$V_{\rm large}/V_{\rm small}=\eta$。
求新共同温度 $T(\eta)$，
假设各腔初始气体物质的量相同、
过程整体可逆。
\exsolution
绝热是指整体与环境没有热交换，
但外力推动活塞仍可向整个系统做功。
因活塞导热，两腔全过程温度在准静态极限下相同，
且可逆过程的整个系统熵守恒。
设两腔各有 $\nu$ 摩尔理想气体，
初始各容积为 $V_0$。
末态总体积固定 $2V_0$，且体积比为 $\eta$，
所以
\[
V_1=\frac{2\eta}{1+\eta}V_0,\quad
V_2=\frac2{1+\eta}V_0.
\]
对两腔分别使用理想气体熵变化，
\[
\Delta S_i=\nu C_V\ln\frac T{T_0}
+\nu R\ln\frac{V_i}{V_0}.
\]
总熵变化为零：
\[
2\nu C_V\ln\frac T{T_0}
+\nu R\ln\frac{V_1V_2}{V_0^2}=0.
\]
因 $R/C_V=\gamma-1$，
\[
\ln\frac T{T_0}
=\frac{\gamma-1}{2}
\ln\frac{(1+\eta)^2}{4\eta}.
\]
因此
\[
\boxed{T=T_0
\left[\frac{(1+\eta)^2}{4\eta}\right]^{(\gamma-1)/2}}.
\]
\excheck 仅靠“总体绝热”不能推出 $U$ 守恒，
因为外界移动活塞\emph{对整个系统做功}；
使用可逆等熵条件才得到本题结果。

\thermoenergyproblem{42}{67}
充有氦气的绝热容器温度为
$T=1000\,\mathrm K$，
气体通过小孔进入真空。
若容器内流速远小于出口速率，
忽略黏性，求出流速率 $v$。
\exsolution
对准稳态绝热无黏流动，
单位质量的总焓和动能守恒：
\[
h_0+\frac12v_0^2=h+\frac12v^2.
\]
题设容器内速度可忽略 $v_0\simeq0$，
并将出流气体对外静态焓的余量在理想真空出口
近似全部转为动能。
理想气体比焓为
$h=C_{p,\rm mol}T/M$，
故
\[
\frac12v^2\simeq\frac{C_{p,\rm mol}T}{M},
\quad
C_{p,\rm mol}=\frac{\gamma R}{\gamma-1}.
\]
因此
\[
\boxed{v=\sqrt{
\frac{2\gamma RT}{(\gamma-1)M}}}.
\]
氦是单原子气体，
$\gamma=5/3$、
$M\simeq4.0\times10^{-3}\,\mathrm{kg/mol}$，
\[
\boxed{v\simeq\sqrt{\frac{5RT}{M}}
\simeq3.2\,\mathrm{km/s}\approx3.3\,\mathrm{km/s}}.
\]
\excheck 这是理想绝热能量转换的流速估计，
不描述实际小孔分子逸出分布的全部动力学细节。

\thermoenergyproblem{43}{67}
一摩尔、热容比为 $\gamma$ 的理想气体
沿 $V=a/T$ 过程变化，
其中 $a$ 为正常数。
若温度上升 $\Delta T$，
求气体净吸收的热量 $Q$。
\exsolution
一摩尔理想气体有
$pV=RT$，由
$V=a/T$ 得
$p=RT/V=RT^2/a$。
沿给定过程的体积微分为
$\mathrm dV=-a\,\mathrm dT/T^2$，
所以气体做功微分
\[
\delta A=p\,\mathrm dV
=\frac{RT^2}{a}
\left(-\frac a{T^2}\mathrm dT\right)
=-R\,\mathrm dT.
\]
内能微分为
$\mathrm dU=C_V\,\mathrm dT
=R\,\mathrm dT/(\gamma-1)$。
根据热力学第一定律，
\[
\delta Q=\mathrm dU+\delta A
=R\left(\frac1{\gamma-1}-1\right)\mathrm dT.
\]
积分得
\[
\boxed{Q=\frac{R(2-\gamma)}{\gamma-1}\Delta T}.
\]
\excheck $T$ 上升使体积 $V=a/T$ 减少，
气体做功为负；不能把 $\delta A$ 当成正值。

\thermoenergyproblem{44}{67}
证明：如果理想气体所做功
与其内能的相应增量成正比，
那么它所经历的准静态过程
可写成多方方程
$pV^n=\text{常数}$。
\exsolution
设比例系数为常数 $k$，
即气体做功微分与内能变化满足
\[
\delta A=p\,\mathrm dV=k\,\mathrm dU.
\]
理想气体内能
$\mathrm dU=\nu C_V\,\mathrm dT$，
而状态方程微分
$\nu R\,\mathrm dT=p\,\mathrm dV+V\,\mathrm dp$。
用
$C_V/R=1/(\gamma-1)$：
\[
p\,\mathrm dV
=\frac{k}{\gamma-1}
(p\,\mathrm dV+V\,\mathrm dp).
\]
对于 $k\ne0$，移项得
\[
\frac{\mathrm dp}{p}
=-\left(1-\frac{\gamma-1}{k}\right)
\frac{\mathrm dV}{V}.
\]
令多方指数
$n=1-(\gamma-1)/k$，
积分为
$\ln p+n\ln V=\text{常数}$，
即
\[
\boxed{pV^n=\text{常数},\qquad
n=1-\frac{\gamma-1}{k}}.
\]
\excheck 如果 $k=0$，
则 $\delta A=0$ 对应体积不变的等容过程，
可视为多方指数 $n\to\infty$ 的极限；
证明中的除以 $k$ 步骤不适用于这一退化情形。

\thermoenergyproblem{45}{67}
一摩尔理想气体的绝热指数为 $\gamma$，
沿多方过程 $pV^n=\text{常数}$
变化。
求过程摩尔热容 $C_n$，
并判断何种多方指数 $n$ 下 $C_n<0$。
\exsolution
多方关系微分为
\[
V\,\mathrm dp+n p\,\mathrm dV=0.
\]
一摩尔理想气体状态方程
$RT=pV$，两边求微分得
\[
R\,\mathrm dT=p\,\mathrm dV+V\,\mathrm dp
=(1-n)p\,\mathrm dV.
\]
因此沿此过程
$\delta A=p\,\mathrm dV=
R\,\mathrm dT/(1-n)$，
而内能微分
$\mathrm dU=C_V\,\mathrm dT
=R\,\mathrm dT/(\gamma-1)$。
按 $C_n=\delta Q/\mathrm dT$，
\[
C_n=
\frac{R}{\gamma-1}+\frac{R}{1-n}
=\boxed{\frac{R(n-\gamma)}
{(n-1)(\gamma-1)}}.
\]
对通常 $\gamma>1$ 且 $n\ne1$，
它为负数当且仅当
\[
\boxed{1<n<\gamma}.
\]
\excheck $n=1$ 是等温过程，
$\mathrm dT=0$ 而可发生热量交换，
因而“摩尔热容 $\delta Q/\mathrm dT$”
在这一点不具通常有限定义；
$n=\gamma$ 为绝热过程，$C_n=0$。
% === END 10-thermo-first-law-036-045.tex ===

% === BEGIN 10-thermo-first-law-046-056.tex blob 4961f6186026c9fec22a56ac9dbd8a7817a484fe ===
% Irodov, user-supplied Arihant PDF, section 2.2, original problems 46--56,
% printed pages 67--69 (PDF pp.76--78). One-mole expressions unless stated.
\thermoenergyproblem{46}{67}
氩气在一次多方过程中，体积增大到原来的
$\alpha=4.0$ 倍，同时压强降至原来的
$1/\beta$，其中 $\beta=8.0$。
把氩气视为理想单原子气体，求此过程的摩尔热容。
\exsolution
设多方过程为 $pV^n=\text{常数}$，
则
\[
\frac{p_2}{p_1}=\left(\frac{V_1}{V_2}\right)^n
\Longrightarrow \frac1\beta=\alpha^{-n},\qquad
n=\frac{\ln\beta}{\ln\alpha}
=\frac{\ln8}{\ln4}=\frac32.
\]
由理想气体第一定律，
过程摩尔热容
\[
C=C_V+\frac{R}{1-n}.
\]
单原子气体 $C_V=3R/2$，
故
\[
\boxed{C=\frac{3R}{2}+\frac{R}{1-3/2}
=-\frac R2\simeq-4.16\,\mathrm{J/(mol\,K)}}.
\]
\excheck 负热容是\emph{此多方路径的热量变化率}
$\delta Q/\mathrm dT<0$，不是单原子气体
平衡定容热容变成负值。

\thermoenergyproblem{47}{68}
一摩尔氩气沿多方指数 $n=1.50$ 的过程膨胀，
温度变化 $\Delta T=-26\,\mathrm K$。
求：(a) 气体吸热 $Q$；(b) 气体做功 $A$。
\exsolution
单原子氩气 $C_V=3R/2$，
多方过程的摩尔热容
\[
C=C_V+\frac R{1-n}
=\frac{3R}{2}-2R=-\frac R2.
\]
因此
\[
\boxed{Q=C\Delta T
=\left(-\frac R2\right)(-26\,\mathrm K)
\simeq+108\,\mathrm J}.
\]
理想气体沿多方过程有
$p\,\mathrm dV=R\,\mathrm dT/(1-n)$，
故
\[
\boxed{A=\frac{R\Delta T}{1-n}
=\frac{8.314(-26)}{-0.50}\,\mathrm J
\simeq432\,\mathrm J}.
\]
\excheck 内能变化
$\Delta U=(3R/2)(-26)\simeq-324\,\mathrm J$，
故 $\Delta U+A\simeq108\,\mathrm J=Q$，
与第一定律一致。

\thermoenergyproblem{48}{68}
绝热指数为 $\gamma$ 的理想气体
沿 $p=\alpha V$ 的规律膨胀，
其中 $\alpha$ 为常数，
初体积 $V_0$，末体积 $\eta V_0$。
求：(a) 内能增量；(b) 气体做功；
(c) 该过程的摩尔热容。
\exsolution
设气体含 $\nu$ 摩尔，则
$pV=\nu RT$、$\nu C_V=\nu R/(\gamma-1)$。
沿 $p=\alpha V$ 有
$\nu RT=\alpha V^2$，所以
\[
\Delta U=\nu C_V\Delta T
=\frac{\alpha}{\gamma-1}(V_2^2-V_0^2).
\]
取 $V_2=\eta V_0$，
\[
\boxed{\Delta U=
\frac{\alpha V_0^2(\eta^2-1)}{\gamma-1}}.
\]
机械功
\[
\boxed{A=\int_{V_0}^{\eta V_0}\alpha V\,\mathrm dV
=\frac{\alpha V_0^2}{2}(\eta^2-1)}.
\]
微分状态方程有
$\nu R\,\mathrm dT=2\alpha V\,\mathrm dV$，
故每摩尔过程热容
\[
C=C_V+\frac{p\,\mathrm dV}{\nu\,\mathrm dT}
=\frac R{\gamma-1}+\frac R2
=\boxed{\frac{R(\gamma+1)}{2(\gamma-1)}}.
\]
\excheck 题目未指定摩尔数；
$\Delta U,A$ 是\emph{整团气体的}能量，
而 $C$ 是每摩尔的热容。

\thermoenergyproblem{49}{68}
绝热指数为 $\gamma$ 的理想气体在膨胀中
吸收的热量恰好等于其\emph{内能减少的量}。
求：(a) 过程摩尔热容；(b) $T,V$ 过程方程；
(c) 一摩尔气体从初温 $T_0$ 出发、
体积增大 $\eta$ 倍时对外做功。
\exsolution
题意为 $\delta Q=-\mathrm dU$，
由第一定律
$\delta Q=\mathrm dU+p\,\mathrm dV$
得到
$p\,\mathrm dV=-2\,\mathrm dU=-2\nu C_V\,\mathrm dT$。
于是
\[
\boxed{C=\frac{\delta Q}{\nu\,\mathrm dT}=-C_V
=-\frac R{\gamma-1}}.
\]
由 $pV=\nu RT$，
\[
\frac{\mathrm dT}{T}
=-\frac{R}{2C_V}\frac{\mathrm dV}{V}
=-\frac{\gamma-1}{2}\frac{\mathrm dV}{V}.
\]
积分：
\[
\boxed{TV^{(\gamma-1)/2}=\text{常数}}.
\]
若为一摩尔，末温
$T_2=T_0\eta^{-(\gamma-1)/2}$，
机械功 $A=-2\Delta U$，故
\[
\boxed{A=\frac{2RT_0}{\gamma-1}
\left[1-\eta^{-(\gamma-1)/2}\right]}.
\]
\excheck 膨胀导致 $T_2<T_0$，
因此 $\Delta U<0$，
$\delta Q=-\mathrm dU>0$，即仍吸热。

\thermoenergyproblem{50}{68}
一摩尔绝热指数为 $\gamma$ 的理想气体，
经历过程 $p=aT^\alpha$，$a,\alpha$ 为常数。
求：(a) 温度增加 $\Delta T$ 时气体做功；
(b) 过程摩尔热容与热容为负的条件。
\exsolution
由理想气体状态方程
$V=RT/p=(R/a)T^{1-\alpha}$，
故
\[
\mathrm dV=(R/a)(1-\alpha)T^{-\alpha}\mathrm dT,
\qquad p\,\mathrm dV=R(1-\alpha)\mathrm dT.
\]
因此
\[
\boxed{A=R(1-\alpha)\Delta T}.
\]
再用
$\delta Q=C_V\,\mathrm dT+p\,\mathrm dV$，
\[
\boxed{C=C_V+R(1-\alpha)
=R\left(\frac{\gamma}{\gamma-1}-\alpha\right)}.
\]
当
\[
\boxed{\alpha>\frac{\gamma}{\gamma-1}}
\]
时 $C<0$。
\excheck $\alpha=1$ 是等容路径，
$A=0$；$\alpha=0$ 是等压路径，
$C=C_p$，均可反查。

\thermoenergyproblem{51}{68}
某理想气体的内能和体积沿过程中满足
$U=aV^\alpha$，
其中 $a,\alpha$ 为常数，
气体绝热指数为 $\gamma$。
若内能增量为 $\Delta U$，求：(a) 气体做功
和传递给气体的热量；(b) 过程摩尔热容。
\exsolution
理想气体满足
$pV=\nu RT=(\gamma-1)U$，
所以
\[
p=\frac{(\gamma-1)U}{V}
=(\gamma-1)aV^{\alpha-1}.
\]
对于 $\alpha\ne0$，
\[
A=\int p\,\mathrm dV
=\frac{\gamma-1}{\alpha}
\,a(V_2^\alpha-V_1^\alpha)
=\boxed{\frac{\gamma-1}{\alpha}\Delta U}.
\]
由第一定律
\[
\boxed{Q=\Delta U+A
=\frac{\alpha+\gamma-1}{\alpha}\Delta U}.
\]
因 $\mathrm dU=\nu C_V\,\mathrm dT$，
过程摩尔热容为
\[
\boxed{C=C_V\frac{\alpha+\gamma-1}{\alpha}
=\frac{R(\alpha+\gamma-1)}
{\alpha(\gamma-1)}}.
\]
\excheck $\alpha=0$ 时沿路径内能不变、
温度不变，此题以非零 $\Delta U$ 给出的
比式不能直接使用，需单独讨论等温极限。

\thermoenergyproblem{52}{68}
某理想气体的定容摩尔热容为 $C_V$。
在下列过程中，求过程摩尔热容随体积 $V$ 的函数：
(a) $T=T_0e^{\alpha V}$；
(b) $p=p_0e^{\alpha V}$。
\exsolution
对一摩尔理想气体，
第一定律给出
\[
C=\frac{\delta Q}{\mathrm dT}
=C_V+p\frac{\mathrm dV}{\mathrm dT}
=C_V+\frac{RT/V}{\mathrm dT/\mathrm dV}.
\]
(a) $T=T_0e^{\alpha V}$，
因此 $\mathrm dT/\mathrm dV=\alpha T$，
\[
\boxed{C_a(V)=C_V+\frac R{\alpha V}}.
\]
(b) $p=p_0e^{\alpha V}$，
状态方程 $T=pV/R$，
故
$\mathrm dT/\mathrm dV=(p/R)(1+\alpha V)$。
代入：
\[
\boxed{C_b(V)=C_V+\frac R{1+\alpha V}}.
\]
\excheck 在 $T'(V)=0$ 的特殊位置，
用 $\delta Q/\mathrm dT$ 定义的过程热容会发散；
与 $C_V$ 常数并不矛盾。

\thermoenergyproblem{53}{68}
一摩尔理想气体的绝热指数为 $\gamma$，
沿 $p=p_0+\alpha/V$ 变化，
其中 $p_0,\alpha>0$。
求：(a) 过程摩尔热容 $C(V)$；
(b) 当 $V_1\to V_2$ 时，
内能增量、气体做功和吸热量。
\exsolution
状态方程给出
$RT=pV=p_0V+\alpha$，
所以 $\mathrm dT=(p_0/R)\,\mathrm dV$，
$\mathrm dV/\mathrm dT=R/p_0$。
于是
\[
\boxed{C(V)=C_V+\frac{Rp}{p_0}
=C_p+\frac{R\alpha}{p_0V}
=\frac{\gamma R}{\gamma-1}+\frac{R\alpha}{p_0V}}.
\]
积分内能变化：
\[
\boxed{\Delta U=\frac{p_0(V_2-V_1)}{\gamma-1}}.
\]
气体做功
\[
\boxed{A=\int_{V_1}^{V_2}
\left(p_0+\frac{\alpha}{V}\right)\mathrm dV
=p_0(V_2-V_1)+\alpha\ln\frac{V_2}{V_1}}.
\]
所以
\[
\boxed{Q=\frac{\gamma p_0}{\gamma-1}(V_2-V_1)
+\alpha\ln\frac{V_2}{V_1}}.
\]
\excheck $p_0V$、$\alpha$ 均有能量量纲，
故每一项都为焦耳。

\thermoenergyproblem{54}{69}
一摩尔理想气体的定压摩尔热容为 $C_p$，
过程满足 $T=T_0+\alpha V$。
求：(a) 过程摩尔热容；
(b) 体积由 $V_1$ 增至 $V_2$ 的吸热量。
\exsolution
由状态方程
$p=RT/V=R(T_0/V+\alpha)$，
且 $\mathrm dT/\mathrm dV=\alpha$。
因此
\[
C=C_V+p\,\frac{\mathrm dV}{\mathrm dT}
=C_p+\frac{RT_0}{\alpha V},
\quad
\boxed{C(V)=C_p+\frac{RT_0}{\alpha V}}.
\]
沿过程吸热
\[
Q=\int_{V_1}^{V_2}C(V)\,\mathrm dT
=\alpha\int_{V_1}^{V_2}C(V)\,\mathrm dV.
\]
计算：
\[
\boxed{Q=C_p\alpha(V_2-V_1)+
RT_0\ln\frac{V_2}{V_1}}.
\]
\excheck 若 $\alpha=0$，
过程为等温且不能把 $C(V)$ 的表达式
当成有限量；但吸热的积分形式仍有极限
$Q=RT_0\ln(V_2/V_1)$。

\thermoenergyproblem{55}{69}
求一摩尔理想气体满足下列过程热容规律时，
$T,V$ 之间的过程方程：
(a) $C=C_V+\alpha T$；
(b) $C=C_V+\beta V$；
(c) $C=C_V+a p$。
\exsolution
对一摩尔理想气体，
$\delta Q=C_V\,\mathrm dT+p\,\mathrm dV$
及 $p=RT/V$ 给出
\[
(C-C_V)\,\mathrm dT
=\frac{RT}{V}\,\mathrm dV.
\]
(a) 代入 $C-C_V=\alpha T$，
\[
\alpha\,\mathrm dT=\frac R V\,\mathrm dV
\Longrightarrow
\boxed{V=V_0e^{\alpha(T-T_0)/R}}.
\]
(b) 若 $C-C_V=\beta V$，
\[
\beta\,\mathrm dT=\frac{RT}{V^2}\,\mathrm dV
\quad\Longrightarrow\quad
\frac{\mathrm dT}{T}=\frac R\beta\frac{\mathrm dV}{V^2}.
\]
积分得到
$\ln T=-R/(\beta V)+\text{常数}$，即
\[
\boxed{T\exp\left(\frac R{\beta V}\right)=\text{常数}}.
\]
(c) 若 $C-C_V=ap$，则
$ap\,\mathrm dT=p\,\mathrm dV$，消去 $p$：
\[
\boxed{V-aT=\text{常数}}.
\]
\excheck (b) 的关键是先除以 $T$ 再积分；
否则会错误得到 $T\propto1/V$。
所得式与原书第55(b)题答案一致。

\thermoenergyproblem{56}{69}
一摩尔理想气体的绝热指数为 $\gamma$，
某过程热容随温度满足 $C=\alpha/T$。
若气体从 $T_0$ 加热到 $\eta T_0$，
求：(a) 气体做功；(b) 以 $p,V$ 表示的过程方程。
\exsolution
由
$\delta Q=C\,\mathrm dT=(\alpha/T)\,\mathrm dT$，
$\mathrm dU=C_V\,\mathrm dT$，
第一定律给
\[
\delta A=p\,\mathrm dV
=\left(\frac{\alpha}{T}-C_V\right)\mathrm dT.
\]
因此
\[
\boxed{A=\alpha\ln\eta-C_VT_0(\eta-1)
=\alpha\ln\eta-\frac{RT_0(\eta-1)}{\gamma-1}}.
\]
另一方面 $p=RT/V$，
\[
R\,\frac{\mathrm dV}{V}
=\left(\frac{\alpha}{T^2}-\frac{C_V}{T}\right)\mathrm dT.
\]
积分得
\[
R\ln V=-\frac\alpha T-C_V\ln T+\text{常数}.
\]
所以
$V T^{1/(\gamma-1)}e^{\alpha/(RT)}=\text{常数}$。
代入 $T=pV/R$ 并对两边取
$\gamma-1$ 次方：
\[
\boxed{pV^\gamma
\exp\left(\frac{(\gamma-1)\alpha}{pV}\right)
=\text{常数}}.
\]
\excheck 若 $\alpha=0$，
过程为绝热，末式退化为
$pV^\gamma=\text{常数}$，
与热力学第一定律一致。
% === END 10-thermo-first-law-046-056.tex ===

% === BEGIN 10-thermo-first-law-057-061.tex blob 9af920a33b7c8a90d15c2837db65f12202090b78 ===
% Irodov, source PDF 78, printed p.69; 2.2 ends at original problem 61.
% One mol van der Waals equation: p = RT/(V-b)-a/V^2.
% For nu mol, U=nu*C_V*T-a*nu^2/V under constant C_V approximation.
\thermoenergyproblem{57}{69}
求一摩尔范德瓦耳斯气体在温度 $T$ 不变时，
由体积 $V_1$ 可逆等温膨胀至 $V_2$ 所做的功。
假定气体遵循
$(p+a/V^2)(V-b)=RT$。
\exsolution
将范德瓦耳斯方程解成压强形式：
\[
p(V,T)=\frac{RT}{V-b}-\frac a{V^2}.
\]
气体在可逆准静态膨胀中所做功为
\[
A=\int_{V_1}^{V_2}p(V,T)\,\mathrm dV
=RT\int_{V_1}^{V_2}\frac{\mathrm dV}{V-b}
-a\int_{V_1}^{V_2}\frac{\mathrm dV}{V^2}.
\]
分别积分，
\[
\boxed{A=
RT\ln\frac{V_2-b}{V_1-b}
+a\left(\frac1{V_2}-\frac1{V_1}\right)}.
\]
\excheck 要求 $V_1,V_2>b$；
$a\to0,b\to0$ 时恢复理想气体的
$A=RT\ln(V_2/V_1)$。

\thermoenergyproblem{58}{69}
一摩尔氧气在温度
$T=280\,\mathrm K$ 不变时，
由体积 $V_1=1.00\,\mathrm L$
膨胀至 $V_2=5.0\,\mathrm L$。
把氧气视为范德瓦耳斯气体，求
(a) 内能变化 $\Delta U$；
(b) 气体吸热 $Q$。
\exsolution
在常定容热容的范德瓦耳斯模型中，
一摩尔内能可写成
\[
U(T,V)=C_VT-\frac aV+\text{常数}.
\]
等温过程中 $\Delta T=0$，
所以
\[
\boxed{\Delta U=a\left(\frac1{V_1}-\frac1{V_2}\right)}.
\]
另一方面第57题给气体做功
\[
A=RT\ln\frac{V_2-b}{V_1-b}
+a\left(\frac1{V_2}-\frac1{V_1}\right).
\]
由第一定律 $Q=\Delta U+A$，
两项中含 $a$ 的部分相消：
\[
\boxed{Q=RT\ln\frac{V_2-b}{V_1-b}}.
\]
取氧气常用范德瓦耳斯参数
$a\simeq1.36\,\mathrm{atm\,L^2/mol^2}$、
$b\simeq0.0318\,\mathrm{L/mol}$：
\[
\Delta U\simeq1.09\,\mathrm{atm\,L}
\simeq\boxed{0.11\,\mathrm{kJ}},
\quad
Q\simeq(8.314)(280)
\ln\frac{5-0.0318}{1-0.0318}
\simeq\boxed{3.8\,\mathrm{kJ}}.
\]
\excheck 与理想气体等温过程的 $\Delta U=0$ 不同：
分子间吸引项 $-a/V$ 导致范德瓦耳斯内能
还依赖体积。

\thermoenergyproblem{59}{69}
范德瓦耳斯气体满足
$(p+a/V^2)(V-b)=RT$（一摩尔）。
在定容摩尔热容 $C_V$ 近似为常数的条件下，求
(a) 可逆绝热曲线在 $T,V$ 坐标中的方程；
(b) 热容差 $C_p-C_V$ 随 $T,V$ 的函数。
\exsolution
由范德瓦耳斯内能
$U=C_VT-a/V$，
有
$\mathrm dU=C_V\,\mathrm dT+(a/V^2)\mathrm dV$。
可逆绝热 $\delta Q=0$，
第一定律给
\[
0=C_V\,\mathrm dT+
\left(\frac a{V^2}+p\right)\mathrm dV
=C_V\,\mathrm dT+\frac{RT}{V-b}\mathrm dV.
\]
分离变量积分：
\[
C_V\frac{\mathrm dT}{T}
=-R\frac{\mathrm dV}{V-b}
\quad\Longrightarrow\quad
\boxed{T(V-b)^{R/C_V}=\text{常数}}.
\]
(b) 使用一般热容差恒等式
\[
C_p-C_V=
-T\frac{(\partial p/\partial T)_V^2}
{(\partial p/\partial V)_T}.
\]
由于
\[
\left(\frac{\partial p}{\partial T}\right)_V=\frac R{V-b},
\qquad
\left(\frac{\partial p}{\partial V}\right)_T
=-\frac{RT}{(V-b)^2}+\frac{2a}{V^3},
\]
得
\[
\boxed{C_p-C_V=
\frac R{1-\dfrac{2a(V-b)^2}{RTV^3}}}.
\]
\excheck 在稀薄气体极限
$a,b\to0$，两结果退化为
$TV^{R/C_V}=\text{常数}$ 和 $C_p-C_V=R$。
热容差表达式仅在稳定的单相分支有物理意义。

\thermoenergyproblem{60}{69}
两个绝热刚性容器由阀门相通：
第一容器体积 $V_1=10\,\mathrm L$，
内有 $\nu=2.5$ mol 二氧化碳；
第二容器体积 $V_2=100\,\mathrm L$，
原来是真空。打开阀门后气体自由膨胀，
假设遵从范德瓦耳斯方程，
求气体的温度变化 $\Delta T$。
\exsolution
两容器总系统绝热且刚性，
向真空自由膨胀不对外做功，
因此总内能守恒：
\[
\Delta U=0.
\]
对 $\nu$ 摩尔的范德瓦耳斯模型，
内能为
\[
U=\nu C_VT-\frac{a\nu^2}{V}+\text{常数}.
\]
最终容积是\emph{两容器之和}
$V_f=V_1+V_2$。
故
\[
0=\nu C_V(T_f-T_i)
-a\nu^2\left(\frac1{V_f}-\frac1{V_1}\right).
\]
因此
\[
\boxed{\Delta T=T_f-T_i=
\frac{a\nu}{C_V}
\left(\frac1{V_1+V_2}-\frac1{V_1}\right)<0}.
\]
采用二氧化碳参数
$a\simeq3.59\,\mathrm{atm\,L^2/mol^2}$，
并与原书此温区热容假设一致地取
$\gamma\simeq1.30$、$C_V=R/(\gamma-1)$，
则
\[
\boxed{\Delta T\simeq-3.0\,\mathrm K}.
\]
\excheck 理想气体向真空自由膨胀温度不变，
但范德瓦耳斯吸引作用导致内能中的
$-a\nu^2/V$ 随体积改变，温度因而降低。
若采用\emph{不同温区的刚性线性分子近似}
$C_V=5R/2$，则会得约 $-4\,\mathrm K$；
因此数值取决于二氧化碳的热容模型，
解析式不依赖这一近似选择。

\thermoenergyproblem{61}{69}
$\nu=3.0$ mol 二氧化碳
在\emph{向真空自由膨胀}时，
由体积 $V_1=5.0\,\mathrm L$
变为 $V_2=10\,\mathrm L$，
欲保持气体温度不变，
需要向气体传入多少热量？
按范德瓦耳斯气体处理。
\exsolution
气体向真空自由膨胀，
对外做功 $A=0$。
范德瓦耳斯总内能
$U=\nu C_VT-a\nu^2/V$，
在 $\Delta T=0$ 时：
\[
\Delta U=
a\nu^2\left(\frac1{V_1}-\frac1{V_2}\right).
\]
由第一定律
$Q=\Delta U+A=\Delta U$，得到
\[
\boxed{Q=a\nu^2\frac{V_2-V_1}{V_1V_2}}.
\]
用 $a\simeq3.59\,\mathrm{atm\,L^2/mol^2}$，
\[
Q\simeq3.59(3.0)^2
\left(\frac15-\frac1{10}\right)
\,\mathrm{atm\,L}
=3.23\,\mathrm{atm\,L}
\simeq\boxed{0.33\,\mathrm{kJ}}.
\]
\excheck \emph{向真空}的自由膨胀与
第57题\emph{可逆等温}膨胀不同。
在自由膨胀中没有准静态 $p\,\mathrm dV$
外功，不能套用第57题的做功积分。
% === END 10-thermo-first-law-057-061.tex ===
